\section{Discussion}
\label{sec:discussion}

At $p=\alpha=1/2$ the numerical evidence supports a discontinuous
active--absorbing transition in a two-dimensional binary rule with
exactly two absorbing configurations. The central observation is
finite, spatially distributed and dynamically sustained activity
across the measured front-reversal region. Prepared fronts reverse
direction while retaining an active interior; bulk decay on the
lower-$q$ side proceeds through absorbing-domain growth; disks and holes
exchange their growth signs on opposite flanks.
The directional crossing differences and their window dependence
limit the resolution of a unique boundary, but do not remove this
combined evidence for phase competition.

The novelty is this joint quantitative test of activity and invasion.
Earlier work reported discontinuous two-state behavior~\cite{rodrigues2015},
and ecological alliance models exhibit related fronts and active
layers~\cite{szabo2004,guisoni2013}. The binary construction isolates
the phenomenon from internal dynamics of an absorbing alliance.
It also distinguishes a discontinuous loss of local activity from
the magnetization jump at a GV transition.

The finite-time closure identifies a specific microscopic contribution:
five or more ancestral lineages generate the noncubic reaction terms
needed for bistability. Its exact identities concern the unperturbed
voter law at a prescribed cutoff. Matching that law to the stochastic
lattice process in two dimensions, especially near cancellation of the
leading drift, remains an approximation. The five-cell comparison
makes the distinction concrete (Supplemental Sec.~\ref{S-app:molofsky}):
a one-sided extrapolation locates the onset of coherent expansion,
but interior decay prevents resolving coherent retreat with a persistent
active interior.
Early seed growth, finite-window bulk stationarity and linear stability
are different quantities.

A predictive spatial theory would need to account for front motion,
curvature and the nucleation rate as well as local drift.
For wetting, it would also need the interaction between the two active
fronts bounding a film. The measured adsorption increase does not fix
that interaction or an asymptotic exponent.
Finally, the near-voter measurements demonstrate why agreement across
sizes at one horizon is insufficient to establish stationarity.
Whether the discontinuity persists for every fixed $p<1$, or ends
before the voter point, is not determined here.
